%! Graphing Quadratic Functions & Transformations — Unit 2, Lesson 2.1 — Homework (Answer Key)
% Mirrors ../homework/main.tex EXACTLY: every \blank{} becomes an \ans{}, every \writelines{n} is
% answered with exactly n \ansline{}s, and every work block is byte-identical, so the two files
% come out the same number of pages. No teacher notes here — they live in the plan.
% GRADUAL RELEASE (LESSON_SHAPE.md §1): this page IS the lesson's individual practice. It is
% scored, it is the last component in the packet, and it is STARTED IN CLASS in the last ten
% minutes of the period, alone and silent, then finished at home. The teacher may substitute a
% DeltaMath set on a given day; that replaces this page and strikes its score cell on the cover.
% Due the first class after two study halls — never "next class".
% TWO PAGES IS A CEILING at 12pt: two boxes — items 1-3 on page 1, items 4-6 and the spiralbox
% on page 2 — so no item breaks across a page. NO EXTENSION BOX (retired with this shape).
% Six items span the standard: the core procedure off a RULE (1), the deliberate CONTRAST PAIR
% carrying the target misconception (2), the same skill off a GRAPH and run backwards (3), the
% SPECIAL CASE and its boundary (4), a MODEL in a fresh context with a work block and an
% interpret-the-answer follow-up (5), and the SOL-STYLE MULTIPLE CHOICE formative check (6).
% NOTHING IS FACTORED OR SOLVED — the factored form in item 4 is GIVEN and read.
% Every work block and every \writelines{n} is matched byte-for-byte / n-for-n in the key.
\documentclass[12pt]{article}
\usepackage{algebra2-article}
\usepackage{algebra2-key}   % pulls in -boxes; do NOT also load -boxes
% Coordinate grid: {xmin}{xmax}{ymin}{ymax}
\newcommand{\qgrid}[4]{%
  \draw[step=1, linegray!40, very thin] (#1,#3) grid (#2,#4);
  \draw[->, semithick, charcoal] (#1,0)--(#2,0) node[below right, font=\scriptsize]{$x$};
  \draw[->, semithick, charcoal] (0,#3)--(0,#4) node[left, font=\scriptsize]{$y$};}
\newcommand{\fdot}[3]{\fill[#1] (#2,#3) circle (3.0pt);}

\begin{document}

\pageheader{Unit 2, Lesson 2.1}{Homework --- Answer Key}
% NAMESTRIP: no \namedateperiod here — mirrors the blank.

\vspace{0.10in}

\boxguard[8]
\begin{remindbox}
\small
\textbf{This is your graded homework.} Your packet with completed homework is due the first
class after two study halls. I will announce the due date in class and on TurtleNet. Remember:
\textbf{read the form you are given.} Vertex form $a(x-h)^2+k$ hands you the vertex $(h,k)$ and
the axis $x=h$ --- \emph{flip the sign inside the bracket}, so $(x-1)^2$ moves the parent
\emph{right} $1$ while $x^2-1$ moves it \emph{down} $1$.
\end{remindbox}

\vspace{0.10in}

\boxguard[14]
\begin{notesbox}{Practice}
\small
\begin{enumerate}[leftmargin=1.9em, itemsep=8pt]
  \item \textbf{Read the vertex form.} For each rule, give the vertex, the axis of symmetry,
        which way it opens, and whether it is narrower or wider than the parent $y=x^2$.
    \\[3pt]
    \vspace{-2pt}
    \begin{center}
    \renewcommand{\arraystretch}{1.35}
    \begin{tabularx}{\linewidth}{@{} l Y Y Y Y @{}}
    \rowcolor{sky}
    \textbf{Rule} & \textbf{Vertex} & \textbf{Axis} & \textbf{Opens} & \textbf{vs.\ parent} \\
    \hline
    (a) $y=4(x-3)^2+2$ & \ans{$(3,2)$} & \ans{$x=3$} & \ans{up} & \ans{narrower} \\
    (b) $y=-\tfrac12(x+1)^2-4$ & \ans{$(-1,-4)$} & \ans{$x=-1$} & \ans{down} & \ans{wider} \\
    (c) $y=-(x-5)^2+7$ & \ans{$(5,7)$} & \ans{$x=5$} & \ans{down} & \ans{same width} \\
    \end{tabularx}
    \end{center}

  \item \textbf{Inside or outside?} Two rules, one character apart: $y=(x+3)^2$ and $y=x^2+3$.
    \vspace{-2pt}
    \begin{center}
    \renewcommand{\arraystretch}{1.3}
    \begin{tabularx}{\linewidth}{@{} l Y Y Y | X @{}}
    \rowcolor{sky}
    \textbf{$x$} & $-3$ & $0$ & $3$ & \textbf{Vertex} \\
    \hline
    \textbf{$y=(x+3)^2$} & \ans{$0$} & \ans{$9$} & \ans{$36$} & \ans{$(-3,0)$} \\
    \textbf{$y=x^2+3$} & \ans{$12$} & \ans{$3$} & \ans{$12$} & \ans{$(0,3)$} \\
    \end{tabularx}
    \end{center}
    (a) One of them moved \emph{sideways} and one moved \emph{up}. Which is which?
        \ans{$(x+3)^2$ moved left $3$; $x^2+3$ moved up $3$} \\[3pt]
    (b) Why does adding inside the bracket move the graph the \emph{other} way? \ansline{The $+3$ is \emph{inside} the bracket, so it changes the input before squaring.}
    \ansline{To get the old output you now need an $x$ that is $3$ smaller --- so the graph slides left.}

  \item \textbf{Read the graph, then write its rule.} The vertex and one other point are marked.
    \\[2pt]
    \begin{minipage}[c]{0.58\linewidth}
    Vertex \ans{$(-2,1)$}; \; axis of symmetry \ans{$x=-2$} \\[3pt]
    Opens \ans{up}; \; $y$-intercept \ans{$(0,5)$} \\[3pt]
    Now write its rule in \textbf{vertex form} (the marked point shows $a=1$): \\[3pt]
    $y=$ \ans{$(x+2)^2+1$}
    \end{minipage}\hfill
    \begin{minipage}[c]{0.38\linewidth}\centering
    \begin{tikzpicture}[scale=0.34]
      \qgrid{-6.5}{2.5}{-1.5}{8.5}
      \begin{scope}
        \clip (-6.5,-1.5) rectangle (2.5,8.5);
        \draw[very thick, forest, domain=-4.8:0.8, samples=70, smooth, variable=\x]
          plot (\x,{(\x+2)^2+1});
      \end{scope}
      \fdot{forest}{-2}{1} \fdot{goldacc}{0}{5}
      \node[below, font=\tiny, charcoal] at (-2,0.7) {$(-2,1)$};
      \node[right, font=\tiny, charcoal] at (0.2,5) {$(0,5)$};
    \end{tikzpicture}
    \end{minipage}
\end{enumerate}
\end{notesbox}

\newpage
\begin{notesbox}{Practice, continued}
\small
\begin{enumerate}[leftmargin=1.9em, itemsep=8pt, start=4]
  \item \textbf{Intercept form, and one special case.} You are \emph{given} the factored rule
        $y=2(x-1)(x+5)$. Do not factor or solve anything --- just read it.
    \begin{enumerate}[label=(\alph*), leftmargin=1.9em, itemsep=4pt]
      \item The zeros are $x=$ \ans{$1$} and $x=$ \ans{$-5$}, so the axis of symmetry
            sits at their midpoint, $x=$ \ans{$\frac{1+(-5)}{2}=-2$}
      \item Substitute that $x$ back in to get the vertex.
        \begin{work}
          y &= 2(-2-1)(-2+5) \\
            &= 2(-3)(3) = -18
        \end{work}
        So the vertex is \ans{$(-2,-18)$}, and it is a \ans{minimum} (maximum / minimum).
      \item \textbf{The special case.} Suppose the two zeros are the \emph{same}, as in
            $y=a(x-3)^2$. How many $x$-intercepts does the parabola have then, and where does the
            vertex sit? \ans{just one, at $x=3$ --- the vertex sits on the $x$-axis}
    \end{enumerate}

  \item \textbf{Model it.} A stone arch over a walkway has height $h(x)=-\tfrac14(x-8)^2+16$
        feet, where $x$ is the distance in feet from the left end of the walkway.
    \begin{enumerate}[label=(\alph*), leftmargin=1.9em, itemsep=4pt]
      \item Straight out of the vertex form: the peak is \ans{$16$} feet high, and it is
            \ans{$8$} feet from the left end.
      \item The arch meets the ground where $h(x)=0$, at $x=0$ and $x=16$. So it is
            \ans{$16$} feet wide at ground level, and the peak sits \ans{halfway across, at $x=8$}.
      \item Find $h(4)$.
        \begin{work}
          h(4) &= -\tfrac14(4-8)^2 + 16 \\
               &= -4 + 16 = 12
        \end{work}
        $h(4)$ means: \ans{the arch is $12$ ft high $4$ ft from the left end}
    \end{enumerate}

  \item \textbf{Multiple choice (SOL-style).} The graph of $y=(x+6)^2$ is the graph of the parent
        $y=x^2$ shifted
    \begin{enumerate}[label=\Alph*., leftmargin=2.2em, itemsep=1pt]
      \item $6$ units to the right.
      \item \ans{$6$ units to the left.} \hfill \textcolor{keyred}{\textbf{$\leftarrow$ correct}}
      \item $6$ units up.
      \item $6$ units down.
    \end{enumerate}
    Say in one sentence how you ruled out one of the others. \ansline{$x+6=x-(-6)$, so $h=-6$: the vertex moves to $(-6,0)$, six units left.}
    \ansline{C and D are out: $+6$ is inside the bracket, and only an outside number moves it up or down.}
\end{enumerate}
\end{notesbox}

\vspace{0.06in}

\boxguard[6]
\begin{spiralbox}
\small
\textbf{Coming next --- Lesson 2.2: Factoring Quadratic Expressions.} Today the factored form
$a(x-p)(x-q)$ was handed to you. Next lesson you build it yourself --- turning $x^2-2x-3$ into
$(x+1)(x-3)$ --- and Lesson~2.3 shows what setting it equal to zero buys you.
\end{spiralbox}

\end{document}
